\documentclass[12pt]{article}
\usepackage[T1]{fontenc}
\usepackage[utf8]{inputenc}
\usepackage{lmodern}
\usepackage{textcomp}
\usepackage{enumitem}
\usepackage{geometry}
\geometry{margin=1in}
\begin{document}
\begin{center}
Answer Key - Chemistry - Graduate Level Assessment 19 \
\small (Answers are ordered exactly as in the question paper.)
\end{center}

\section*{Multiple Choice}
\begin{enumerate}[label=\arabic*.]
\item \textbf{Answer:} A
\\ \textit{Options:} A) 1.256 \ensuremath{\times} 10^-46 kg $m^2$ and 9.100 \ensuremath{\times} 10^-2 nm; B) 1.336 \ensuremath{\times} 10^-47 kg $m^2$ and 8.250 \ensuremath{\times} 10^-2 nm; C) 1.336 \ensuremath{\times} 10^-46 kg $m^2$ and 9.196 \ensuremath{\times} 10^-1 nm; D) 1.410 \ensuremath{\times} 10^-47 kg $m^2$ and 9.250 \ensuremath{\times} 10^-2 nm; E) 1.256 \ensuremath{\times} 10^-47 kg $m^2$ and 8.196 \ensuremath{\times} 10^-2 nm
\\ \textit{Note:} The correct answer is derived from the relationship between the rotational spectrum spacing and the moment of inertia, where the line spacing of 41.9 cm^-1 corresponds to the rotational energy levels given by the formula \(E = \frac{h^2}{8\pi^2 I} J(J+1)\), allowing for the calculation of both the moment of inertia and the bond length of HF using its reduced mass.
\item \textbf{Answer:} E
\\ \textit{Options:} A) 4 x 10^-15 J; B) 1.2 x 10^-15 J; C) 5 x 10^-15 J; D) 3.5 x 10^-15 J; E) 2 $10^{-15} \mathrm{~J}$
\\ \textit{Note:} The energy of a photon can be calculated using the equation \(E = \frac{hc}{\lambda}\), where \(h\) is Planck's constant, \(c\) is the speed of light, and \(\lambda\) is the wavelength, so for a wavelength of 100 pm (or \(100 \times 10^{-12}\) m), the resulting energy calculation yields approximately \(2 \times 10^{-15} \mathrm{~J}\).
\item \textbf{Answer:} A
\\ \textit{Options:} A) the protons, electrons and neutrons; B) the electrons and protons; C) the nucleus; D) the protons; E) the neutrons
\\ \textit{Note:} The correct answer is based on the fact that the mass of electrons is negligible compared to protons and neutrons, and when combined, the total mass contribution from all three subatomic particles is still less than 1/1000 of the atom's total mass, as most of the mass comes from the nucleus, which is comprised primarily of protons and neutrons.
\item \textbf{Answer:} A
\\ \textit{Options:} A) -2; B) 1; C) -0.5; D) 0.5; E) 2
\\ \textit{Note:} The electromotive force (emf) of a system is determined by the Nernst equation, which states that when the reaction quotient Q equals the equilibrium constant K, the emf is zero; however, if given a specific value like -2, it suggests that the system is under non-standard conditions where the Gibbs free energy change is negative, indicating a spontaneous reaction.
\item \textbf{Answer:} A
\\ \textit{Options:} A) 0.11 J/g^\ensuremath{\circC}; B) 0.20 J/g^\ensuremath{\circC}; C) 0.13 J/g^\ensuremath{\circC}; D) 0.18 J/g^\ensuremath{\circC}; E) 0.30 J/g^\ensuremath{\circC}
\\ \textit{Note:} The heat capacity of uranium is calculated by dividing its atomic heat (26 J/mole) by its molar mass (238 g/mole), resulting in a specific heat capacity of 0.11 J/g°C.
\end{enumerate}

\section*{True / False}
\begin{enumerate}[label=\arabic*.]
\setcounter{enumi}{5}
\item \textbf{Answer:} False
\\ \textit{Options:} A) True; B) False
\\ \textit{Note:} The oxidation number of chlorine in the perchlorate ion, ClO 4-, is +7, not +5, because the overall charge of the ion is -1 and the four oxygen atoms each have an oxidation number of -2.
\item \textbf{Answer:} False
\\ \textit{Options:} A) True; B) False
\\ \textit{Note:} The statement is false because in liquid ammonia, the strongest base is the amide ion (NH 2-) rather than ammonia itself, as it can deprotonate ammonia more effectively.
\item \textbf{Answer:} False
\\ \textit{Options:} A) True; B) False
\\ \textit{Note:} The statement is false because an atom of 52 Cr (chromium) has 24 protons, 28 neutrons, and 24 electrons, as its atomic number is 24 and the mass number is 52, which indicates that it contains 28 neutrons (52 - 24 = 28).
\item \textbf{Answer:} True
\\ \textit{Options:} A) True; B) False
\\ \textit{Note:} The statement is true because SnF\_4, PH\_3, and Si\_2 O\_5 are the correct empirical formulas formed from the respective elements, with tin (Sn) forming a tetravalent compound with fluorine (F), phosphorus (P) forming a trivalent hydride with hydrogen (H), and silicon (Si) exhibiting a higher oxidation state in its oxide with oxygen (O).
\item \textbf{Answer:} False
\\ \textit{Options:} A) True; B) False
\\ \textit{Note:} A T$_{2}$ value of 15 ms corresponds to a linewidth of approximately 10.9 Hz, calculated using the relation linewidth (Hz) = 1/(2πT 2), which indicates that the statement is false.
\end{enumerate}

\section*{Long Form Response}
\begin{enumerate}[label=\arabic*.]
\setcounter{enumi}{10}
\item \textbf{Answer:} The 1 H NMR spectrum of chloroform (CHCl 3) appears as a singlet due to the presence of both 35 Cl and 37 Cl isotopes, each of which has an electric quadrupole moment. This quadrupole moment leads to the interaction of the chlorine nuclei with the magnetic field, causing them to have a non-spherical charge distribution. As a result, the nuclear spin states of the chlorine isotopes do not couple with the hydrogen nuclei in a way that would create additional splitting in the NMR spectrum. Consequently, the hydrogen atom experiences a uniform magnetic environment, leading to the observation of a singlet peak in the 1 H NMR spectrum. This behavior highlights the importance of isotopic effects and electric quadrupole interactions in NMR spectroscopy.
\\ \textit{Note:} The 1 H NMR spectrum of CHCl 3 appears as a singlet because the electric quadrupole moments of both 35 Cl and 37 Cl prevent the chlorine nuclei from coupling with the hydrogen nuclei, thus avoiding any additional splitting in the spectrum.
\item \textbf{Answer:} To calculate the relative abundance of the isotopes ^210 X and ^212 X for the unknown element X, we set up a system of equations based on the average atomic weight. Let the abundance of ^210 X be represented as x and that of ^212 X as (1 - x). The equation can be formulated as: 210.197 amu = 209.64 amu * x + 211.66 amu * (1 - x). Solving this yields x \ensuremath{\approx} 0.7242, indicating that approximately 72.42\% of the element is ^210 X, while 27.58\% is ^212 X. This relative abundance suggests a predominant presence of the more stable isotope, ^210 X, which may imply that the element has a tendency toward stability, as the higher abundance of stable isotopes often correlates with lower radioactivity and greater longevity in nature.
\\ \textit{Note:} correct because it uses the weighted average formula to accurately determine the relative abundances of the isotopes based on their masses, leading to the conclusion that the majority of the element consists of the more stable isotope ^210 X, which has implications for the element's overall stability.
\end{enumerate}

\vfill
\noindent\textit{End of Answer Key}
\end{document}